\section{Justificación de la tecnología seleccionada}

Se usa PostgreSQL con \texttt{pgvector} y \texttt{btree\_gist}. El caso requiere
relaciones, vigencia temporal, autorización por fila, agregaciones SQL y
búsqueda vectorial. Un solo motor aplica esas reglas en la misma transacción y
evita sincronizar permisos y estados con un índice externo.

\begin{table}[ht]
\centering
\small
\begin{tabularx}{\textwidth}{>{\raggedright\arraybackslash}p{0.22\linewidth}Y Y}
\toprule
Alternativa & Ventaja posible & Motivo para no adoptarla como motor principal \\
\midrule
Base documental & Flexibilidad para versiones y metadatos JSON. & Traslada relaciones, exclusión temporal y autorización uniforme a lógica de aplicación. \\
Clave--valor & Baja latencia para caché o contexto efímero. & No resuelve joins, agregaciones ni trazabilidad histórica como almacenamiento primario. \\
Columnar & Escritura sostenida de eventos y consultas por rango temporal. & La muestra necesita correlacionar auditoría con actores, consultas y respuestas mediante joins. \\
Grafo & Recorridos de relaciones variables. & Las relaciones del caso tienen profundidad conocida y se expresan con claves foráneas. \\
Vectorial dedicada & Mayor volumen de búsquedas por segundo a gran escala. & Exige replicar vigencia y permisos para filtrar antes del top-k, con riesgo de inconsistencia entre motores. \\
\bottomrule
\end{tabularx}
\caption{Alternativas consideradas frente a PostgreSQL.}
\end{table}

\texttt{pgvector} aporta \texttt{VECTOR(32)}, distancia coseno y el índice HNSW.
\texttt{btree\_gist} combina igualdad e intervalos en restricciones de
exclusión. PostgreSQL también ofrece transacciones, claves foráneas, constraint
triggers, roles y RLS. Estas funciones son centrales para el caso, no accesorios
de la búsqueda semántica.
